\chapter{Viability}

\section{What viability is for}

Variation proposes candidates. Viability decides who stays. The
viability slot --- $F : \mathcal{M}_+(X) \to \mathcal{M}_+(X)$, a map
from population measures to population measures --- is the rule that
turns ``here is everyone who was just born or moved'' into ``here is
everyone who is alive at the start of the next tick.''

This is the slot with the most contested content in the lens. Chapter~2
named four traditions for thinking about what life is, and three of
them --- Schr\"odinger's persistence, autopoietic closure, the free
energy principle --- have direct formal viability proposals. RAF set
theory adds a fourth. The book takes them as four distinct formalisms.
Whether they secretly collapse to a common closure operator is a real
question, and an unresolved one: it is Conjecture~C1 in the paper (see
\texttt{RESULTS.md}). We will set out the conjecture and the partial
evidence; we will not pretend it is proven.

The chapter has three parts. First, what closure operators are and why
they matter for viability. Second, the four formalisms in turn. Third,
what to do when the viability slot is hierarchical --- a composite
entity made of sub-entities, each of which has its own filter.

\section{Closure operators, gently}

Three of the four viability formalisms are closure conditions of some
kind. The math primer at the back of the book has a full treatment;
here is the working version.

A \emph{partially ordered set} (or poset) is a set with a binary
relation $\sqsubseteq$ that is reflexive ($A \sqsubseteq A$ for all
$A$), anti-symmetric (if $A \sqsubseteq B$ and $B \sqsubseteq A$ then
$A = B$), and transitive (if $A \sqsubseteq B$ and $B \sqsubseteq C$
then $A \sqsubseteq C$). The canonical example: the power set $2^P$ of
any set $P$, ordered by inclusion $\subseteq$. Bigger sets are
``higher.''

A \emph{closure operator} on a poset is a function $c : 2^P \to 2^P$
satisfying three conditions:

\begin{itemize}[leftmargin=1.5em,itemsep=2pt]
  \item \textbf{Extensivity}: $A \subseteq c(A)$. The closure can only
    grow, never shrink.
  \item \textbf{Monotonicity}: if $A \subseteq B$, then $c(A) \subseteq
    c(B)$. Closure respects inclusion.
  \item \textbf{Idempotence}: $c(c(A)) = c(A)$. Closing twice is the
    same as closing once.
\end{itemize}

A set $A$ is \emph{closed} (under $c$) if $A = c(A)$. The closed sets
form a sub-poset of $2^P$ and have the property that an intersection of
closed sets is closed.

The intuition: $c$ takes a candidate set and ``fills in everything that
follows.'' If you give $c$ a small set, it returns the smallest closed
set containing it. The closed sets are the things that are
\emph{stable} under the implicit rule of $c$.

Several of the viability formalisms in this chapter say, in their own
language, ``an entity passes the filter if its supporting set $S$ is
closed under a relevant closure operator $c$.'' The framework
conjectures (C1) that this is the right way to read all four
formalisms. Here is the chapter.

\section{Markov-blanket viability}

The Friston free energy principle says, roughly, that anything that
persists must admit a Markov blanket. A Markov blanket is a
decomposition of a system's variables into four sets:

\begin{itemize}[leftmargin=1.5em,itemsep=2pt]
  \item \textbf{Internal} states $\mu$: the inside of the entity.
  \item \textbf{Active} states $a$: the entity's actions on the
    outside.
  \item \textbf{Sensory} states $s$: signals the entity reads from the
    outside.
  \item \textbf{External} states $\eta$: everything else.
\end{itemize}

The defining property: $\mu$ is conditionally independent of $\eta$
given $s \cup a$ (which together form the \emph{blanket}). The blanket
sits between the inside and the outside; it screens off direct
dependencies.

Markov-blanket viability is the rule: an entity passes the filter iff
its supporting set $S$ admits this kind of sparse-coupling
decomposition into $(\mu, a, s, \eta)$ with the right conditional
independences, and that decomposition is preserved under the substrate
dynamics $\Phi$.

\subsection*{Why this is a viability rule}

If the decomposition is preserved over time, the system has a stable
inside that is causally insulated from the outside in a particular,
information-theoretic way. The internal states can be interpreted as
parameterising a Bayesian generative model whose free energy they
minimise --- this is the active-inference interpretation. The
generative-model interpretation is consequence; the conservation of
the blanket is the constraint.

\subsection*{Caveats}

Aguilera, Biehl, and Heins have argued that the conditions under which
sparse coupling implies a genuine Markov blanket structure are
stronger than the FEP literature often acknowledges. Their work
identifies cases where sparse coupling at the level of system
equations does not yield blanket structure at the level of statistical
dependencies. Friston's reply sharpens but does not resolve the
dispute.

The lens takes Markov-blanket viability as one of four candidate
formalisms. It does not require an entity to admit a Markov blanket;
it requires the implementer to declare whether they are using this
formalism, and to verify the blanket conditions on their substrate
when they do.

\subsection*{Is Markov-blanket viability a closure?}

This is the case where Conjecture~C1 is genuinely open. The
statistical decomposition into $(\mu, a, s, \eta)$ is not obviously a
closure operator on a poset of supporting sets. One can imagine
reformulations that would make it one --- ``the smallest set
containing the candidate inside that is causally screened from its
outside'' --- but the equivalence is unproven.

\section{Autopoietic closure}

The Maturana--Varela autopoietic closure (Chapter~2) says: a system is
viable iff it is closed under its own constitutive network. The
network's products are exactly the network's inputs. The system makes
itself.

For artificial-life systems, Randall Beer has formalised autopoietic
closure for Conway's Game of Life. The reaction network of the GoL is
the rule mapping a cell's neighbourhood to its next state. An entity's
supporting set is closed under the preimage of this reaction network
if every cell in the set's update depends only on cells in the set
(or in the immediately bordering region whose state is fully determined
by the set). A glider is closed in this sense: its support's update is
fully determined by the support's current state, modulo a
deterministic boundary.

\subsection*{Autopoietic closure as a closure operator}

This case is the cleanest fit for C1. The relation ``site $p$ depends
on the set $S$ for its next-state'' generates a preimage-closure on
the reaction-set lattice; the closed sets are exactly the autopoietic
units. The math primer covers the construction; the upshot is that
this formalism instantiates the closure-operator schema directly.

\subsection*{Applying autopoietic closure}

In practice, identifying autopoietic units of a continuous substrate
(Lenia, Flow-Lenia) is harder than for the GoL. You need:

\begin{itemize}[leftmargin=1.5em,itemsep=2pt]
  \item A way of identifying candidate supporting sets $S$ (typically
    connected components above some threshold).
  \item A specification of the local ``reaction network'' on
    $\Phi$. For Lenia, this is the kernel convolution: site $p$'s next
    state depends on sites within the kernel radius of $p$.
  \item A check that $S$ is closed under this reaction relation: every
    site in $S$'s update depends only on sites in $S$, and every site
    that depends on $S$ is in $S$ (modulo a deterministic boundary
    that the substrate dynamics handles).
\end{itemize}

The cleanest Lenia case is a mass-conservation closure: a Flow-Lenia
creature is autopoietically closed in the sense that mass is conserved
inside its support (it does not import or export mass to the
surrounding field).

\section{RAF set viability}

RAF stands for Reflexively Autocatalytic and Food-generated. The
notion comes from chemistry and the work of Stuart Kauffman, Wim
Hordijk, and Mike Steel. Strip the chemistry away and the condition is
the following.

You have a reaction network --- a directed bipartite graph between
chemical species and reactions. Each reaction has some reactants, some
products, and possibly a catalyst. There is a distinguished \emph{food
set} $F \subseteq$ species (the always-available inputs).

A subset $R' \subseteq$ reactions is a \emph{RAF set} if:

\begin{itemize}[leftmargin=1.5em,itemsep=2pt]
  \item \textbf{Reflexively autocatalytic}: every reaction in $R'$ is
    catalysed by some species that is either in the food set or
    produced by another reaction in $R'$.
  \item \textbf{Food-generated}: every reactant of every reaction in
    $R'$ is either in the food set or produced by another reaction
    in $R'$.
\end{itemize}

So $R'$ is a self-supporting, food-fed sub-network. The condition is
exactly a closure: $R'$ is the smallest set containing itself that is
closed under the relations of mutual catalysis and food-generation.

\subsection*{RAF as a closure}

This case is essentially proved. The maximum-RAF (the largest RAF
sub-network of a reaction network) is exactly the closure of an
appropriately-defined operator on the lattice of reaction subsets.
Hordijk and Steel's polynomial-time algorithm for finding the maxRAF
is a closure-iteration: start with all reactions, repeatedly remove
ones that fail the conditions, halt when nothing more is removable.

The fixed point is the closure.

\subsection*{Using RAF for non-chemical substrates}

RAF was developed for chemical reaction networks, but the structure
applies wherever you have a network of mutually-supporting processes
with some always-available input. For digital substrates, this might
be:

\begin{itemize}[leftmargin=1.5em,itemsep=2pt]
  \item In Tierra or Avida, the ``food set'' is the set of always-
    available CPU instructions; reactions are program transformations;
    autocatalysis is ``the program produces itself.''
  \item In LLM-policy substrates, the ``food set'' is the base model's
    capabilities; reactions are prompt mutations; autocatalysis is
    ``the policy promotes the variations that produce it.''
\end{itemize}

The book does not pursue these extensions, but they are natural and
under-explored.

\section{Minimal Criterion Coevolution}

Minimal Criterion Coevolution (MCC), due to Jonathan Brant and
Kenneth Stanley, is the simplest of the four formalisms. An entity is
viable if it satisfies a minimal qualifying predicate, and the
predicate's content is supplied by another co-evolving population.

A worked example: solvers and mazes. The solver population evolves
under the predicate ``solves at least one maze in the current maze
population.'' The maze population evolves under the predicate ``is
solved by at least one solver in the current solver population.'' Each
side's predicate is supplied by the other; both sides face minimal
qualifying criteria; the result is open-ended-ish coevolution where
the difficulty of mazes and the skill of solvers grow together.

\subsection*{Is MCC a closure?}

This is where C1 is most strained. MCC is stated as a predicate, not
as a closure. Can it be rewritten as one?

One candidate reformulation: take the lattice of (solver, maze) pairs;
say a pair is ``closed'' if the solver in the pair solves the maze in
the pair. The viable sub-populations are unions of pairs from the
closed set. Whether this construction gives a genuine closure operator
(extensive, monotone, idempotent) in the relevant sense is unproven.

For purposes of the lens, treat MCC as predicate-based viability rather
than closure-based viability, and accept that C1 is genuinely open for
this case.

\section{Hierarchical viability}

The four formalisms apply at one level. Real systems are nested:
multicellular organisms have viable cells; ant colonies have viable
ants and viable colonies; LLM agents in cultural-evolution
simulations have viable individuals and viable communities. The lens
handles this via \emph{hierarchical viability}.

Recall from Chapter~4 that entities compose. If $e_1, \ldots, e_k$ are
entities with disjoint supporting sets, their union (with joint scale
projection) is a candidate composite entity $e^* = (\bigcup S_i,
(\pi_1, \ldots, \pi_k))$. The composite entity satisfies the entity
definition iff the joint information-closure condition holds at the
joint scale.

The hierarchical viability filter $F^*$ accepts $e^*$ iff:

\begin{itemize}[leftmargin=1.5em,itemsep=2pt]
  \item each component $e_i$ satisfies its own filter $F_i$, and
  \item the joint configuration satisfies a top-level filter
    $F_{\mathrm{top}}$.
\end{itemize}

In words: each part must be viable, and the whole must additionally be
viable as a whole. This is the obvious move; whether it works
correctly under all four formalisms is the second open conjecture, C2
(see \texttt{RESULTS.md}).

The status:

\begin{itemize}[leftmargin=1.5em,itemsep=2pt]
  \item Beer's autopoietic gliders compose correctly under simple
    unions. C2 looks plausible here.
  \item RAFs do \emph{not} compose under union in general: the union
    of two RAFs is not necessarily a RAF without their joint
    catalyses. This is a partial counterexample for one direction of
    C2.
  \item Markov-blanket viability under composition is open.
  \item MCC under composition is open.
\end{itemize}

The lens accommodates hierarchical viability by allowing the user to
specify $F^*$. It does not yet guarantee that $F^*$ behaves the way
intuition expects across formalisms.

\section{The unification conjecture C1, and why the book leaves it open}

We have surveyed four viability formalisms. Three of them (autopoietic
closure, RAF, and arguably Markov blanket) admit some kind of
closure-operator description. Are they secretly the same notion?

Conjecture~C1: \emph{each of the four viability formalisms corresponds
to a closure operator $c_\bullet$ on a partial order of supporting
sets, in the sense that an entity passes the filter iff its supporting
set is closed under $c_\bullet$.}

A proof of C1 would unify the four formalisms abstractly: viability
would have a single, formal definition, and the four named instances
would be choices of $c_\bullet$. A refutation --- a counterexample for
one of the cases --- would reveal that ``viability'' is genuinely
heterogeneous in the same way that ``life'' is heterogeneous across
biological taxa.

The book \emph{deliberately does not assume C1}. There is a temptation
to ship a unified \texttt{ClosureViability} class in the scaffold and
have the four formalisms be subclasses; this would be a mistake, and
the implementation explicitly avoids it. The four formalisms are
separate protocols. If C1 is later established, the lens can be
refactored without changing user code; if it is refuted, the framework
has lost nothing.

\subsection*{Where C1's evidence stands}

\begin{itemize}[leftmargin=1.5em,itemsep=2pt]
  \item \emph{Autopoietic closure}: explicit preimage closure on the
    reaction-set lattice. C1 holds here.
  \item \emph{RAF}: maxRAF is a closure-iteration fixed point. C1
    holds here.
  \item \emph{Markov blanket}: statistical decomposition is not
    obviously a closure. C1 is open here.
  \item \emph{MCC}: predicate-based. C1 is open here.
\end{itemize}

Two out of four is partial evidence, not proof. We will return to C1
in Chapter~9, where the rest of the open problems live.

\section{The practical consequence}

The four formalisms are \emph{not in general substitutable in code}.
A simulator that swaps Markov-blanket viability for autopoietic
closure must verify that the substitution preserves the predicates
the rest of the system relies on. The system description in the
paper requires the viability formalism to be named explicitly, not
merely as ``$F : \mathcal{M}_+(X) \to \mathcal{M}_+(X)$.''

The cost of this caution is a small amount of extra ceremony when
specifying a system. The benefit is that you do not silently swap
one philosophy of life for another.

\section*{To think about}

\begin{enumerate}
  \item Pick a Lenia creature you find aesthetically interesting.
    Write down what a candidate supporting set $S$ would look like.
    Now pick one of the four viability formalisms and check whether
    your $S$ is closed under it.
  \item C1 is open for two of the four cases. Construct (in prose) a
    plausible counterexample to C1 for Markov-blanket viability:
    something that should pass an FEP-style viability test but is
    \emph{not} closed under any natural closure operator.
  \item Hierarchical viability accepts an entity iff (a) each
    component is viable and (b) the whole is viable as a whole. Why
    is (b) needed? Construct a case where (a) holds but (b) fails.
  \item RAF was developed for chemical networks. Sketch how to apply
    it to Tierra. What is the food set?
  \item The book treats the four formalisms as separate. A unified
    framework would make code simpler. What is the risk of unifying
    too early?
\end{enumerate}
